\section{System Overview}
\label{sec:overview}

%% fig-layers.tex -- Figure 1: the six-layer stack (full width).
%%
%% Read bottom-up: infrastructure carries kernel services; kernel services obey
%% a governance substrate; the substrate governs asset families (the nouns);
%% lifecycle modes (the verbs) act on those nouns; surfaces expose the result.
%% The left rail holds platform services that cut across all six layers rather
%% than occupying one; its dashed outline marks that permeability.
%%
%% Type scale: every label is \figH (8pt) or \figB (7pt) -- see preamble.tex. The
%% content was trimmed to fit that floor rather than the floor lowered to fit the
%% content, because everything that came out is restated in the prose.

\begin{figure}[tbp]
\centering
\ifcbexcalidraw
  \begin{cbwide}\cbdiagram{fig-layers}\end{cbwide}
\else
\begin{cbwide}
\begin{tikzpicture}[x=1cm,y=1cm]

  % ---- geometry ------------------------------------------------------------
  % All box widths derive from one content width, so a band stays flush when a
  % label changes length. Total width stays inside the 7.0in text block.
  \def\CX{2.62}          % content left edge
  \def\CW{12.34}         % content width
  \def\G{0.18}           % gutter between boxes
  \def\LX{15.12}         % layer-label column left edge

  % ---- L6 surfaces (4) -----------------------------------------------------
  \foreach [count=\i from 0] \t in {%
      {\textbf{React SPA}\\{\figB same origin}},
      {\textbf{HTTP API}\\{\figB \statRouteDecls{} routes}},
      {\textbf{\aria copilot}\\{\figB tool loop}},
      {\textbf{Headless}\\{\figB SSE, webhooks}}}
  {
    \pgfmathsetmacro{\w}{(\CW-3*\G)/4}
    \node[surface, anchor=north west, text width=\w cm-8pt, minimum height=0.96cm]
      at ({\CX+\i*(\w+\G)}, 0) {\t};
  }

  % ---- L5 lifecycle modes (6) ---------------------------------------------
  \foreach [count=\i from 0] \t in {%
      {\mode{Author}\\{\figB draft, validate}},
      {\mode{Test}\\{\figB bounded runs}},
      {\mode{Evaluate}\\{\figB score, judge}},
      {\mode{Calibrate}\\{\figB optimize}},
      {\mode{Release}\\{\figB gate, promote}},
      {\mode{Observe}\\{\figB trace, meter}}}
  {
    \pgfmathsetmacro{\w}{(\CW-5*\G)/6}
    \node[control, anchor=north west, text width=\w cm-8pt, minimum height=0.96cm,
          inner xsep=2pt] at ({\CX+\i*(\w+\G)}, -1.34) {\t};
  }

  % ---- L4 asset families (9) ----------------------------------------------
  \foreach [count=\i from 0] \t in {%
      {Prompt}, {Tool}, {Skill}, {MCP}, {Work-\\flow},
      {Know.\\base}, {Test\\set}, {Judge}, {Agent}}
  {
    \pgfmathsetmacro{\w}{(\CW-8*\G)/9}
    \node[asset, anchor=north west, text width=\w cm-8pt, minimum height=0.84cm,
          font=\figB, inner xsep=1.5pt] at ({\CX+\i*(\w+\G)}, -2.68) {\t};
  }

  % ---- L3 governance (5) ---------------------------------------------------
  \foreach [count=\i from 0] \t in {%
      {\textbf{Identity}\\{\figB \statScopes{} RBAC scopes}},
      {\textbf{Exec.\ policy}\\{\figB governed egress}},
      {\textbf{Evidence}\\{\figB gate verdicts}},
      {\textbf{Release}\\{\figB promote, revert}},
      {\textbf{Ledgers}\\{\figB audit, effects}}}
  {
    \pgfmathsetmacro{\w}{(\CW-4*\G)/5}
    \node[govern, anchor=north west, text width=\w cm-8pt, minimum height=0.96cm,
          inner xsep=2pt] at ({\CX+\i*(\w+\G)}, -3.90) {\t};
  }

  % ---- L2 kernel (7) -------------------------------------------------------
  \foreach [count=\i from 0] \t in {%
      {Config \&\\secrets}, {Persis-\\tence}, {Storage}, {Tool\\exec.},
      {Event\\bus}, {Observ-\\ability}, {Provider\\adapters}}
  {
    \pgfmathsetmacro{\w}{(\CW-6*\G)/7}
    \node[control, anchor=north west, text width=\w cm-8pt, minimum height=0.96cm,
          font=\figB, inner xsep=1.5pt] at ({\CX+\i*(\w+\G)}, -5.24) {\t};
  }

  % ---- L1 infrastructure (5) ----------------------------------------------
  \foreach [count=\i from 0] \s/\t in {%
      store/{Starlette\\{\figB ASGI host}},
      store/{PostgreSQL 17\\{\figB pgvector, AGE}},
      store/{Object store\\{\figB S3, local}},
      extern/{\mlflow{} 3.14+\\{\figB registry, traces}},
      async/{Transport\\{\figB NATS, Redis}}}
  {
    \pgfmathsetmacro{\w}{(\CW-4*\G)/5}
    \node[\s, anchor=north west, text width=\w cm-8pt, minimum height=0.96cm,
          inner xsep=2pt] at ({\CX+\i*(\w+\G)}, -6.58) {\t};
  }

  % ---- the "rests on" spine ------------------------------------------------
  \foreach \ya/\yb in {-1.00/-1.30, -2.34/-2.64, -3.56/-3.86,
                       -4.90/-5.20, -6.24/-6.54}
    { \draw[flow, draw=cbMuted] (\CX+\CW/2, \ya) -- (\CX+\CW/2, \yb); }

  % ---- layer labels (right column) ----------------------------------------
  \foreach \y/\num/\name/\gloss in {
      0.00/6/SURFACES/the interface,
     -1.34/5/{LIFECYCLE MODES}/the verbs,
     -2.68/4/{ASSET FAMILIES}/the nouns,
     -3.90/3/GOVERNANCE/the rules,
     -5.24/2/KERNEL/services,
     -6.58/1/INFRASTRUCTURE/the base}
  {
    \node[anchor=north west, align=left, text width=2.66cm, font=\figB,
          text=cbInk, inner sep=0pt] at (\LX, \y-0.08)
         {\hyphenpenalty=10000\exhyphenpenalty=10000
          \textbf{\num}\dt\textbf{\name}\\[-0.1ex]
          {\itshape\color{cbMuted}\gloss}};
  }

  % ---- cross-cutting platform services (left rail) -------------------------
  % Dashed outline: these services are not a layer -- every layer reaches them.
  \node[muted, densely dashed, anchor=north west, text width=1.98cm, align=left,
        font=\figB, minimum height=7.54cm, inner xsep=3pt] at (0.0, 0.0)
    {\textbf{PLATFORM}\\\textbf{SERVICES}\\[0.4ex]
     {\itshape not a layer ---}\\{\itshape every layer}\\{\itshape reaches them}\\[1.0ex]
     Evidence base\\[0.45ex]
     Evaluation\\[0.45ex]
     Calibration\\[0.45ex]
     Capability\\registry\\[0.45ex]
     Integration hub\\[0.45ex]
     Project scoping};

\end{tikzpicture}
\end{cbwide}
\fi
\caption{The \sys stack. Six layers, read bottom-up: \Linfra{} carries
\Lkernel{}; kernel services obey \Lgov{}; the substrate governs \Lfam{} --- the
nouns; \Lmode{} --- the verbs --- act on those nouns; and \Lsurf{} exposes the
result. The dashed left rail holds platform services that cut across all six
layers rather than occupying one. Adjacency in this diagram is \emph{capability
availability, not inheritance}: a family in \Lnum{4} obtains \Lnum{5} behaviour
and \Lnum{3} enforcement only by explicitly wiring them, which is why the
guarantees in \tabref{tab:families} differ per row instead of being uniform. That
distinction is the paper's central architectural claim (\secref{sec:arch}).}
\label{fig:layers}
\figkey
\end{figure}

%% fig-system.tex -- Figure: high-level system architecture and component
%% interaction. One ASGI process, the clients that reach it, and the systems it
%% integrates with but does not own.
%%
%% Routing convention: every long edge runs either in the gutter between the
%% process band and the state column (lane A), in the outer lane to the right of
%% everything (lane B), or in the reserved channel below the band. Nothing
%% crosses a box.
%%
%% Type scale: \figH headings, \figB detail (7pt floor). Row detail lines were cut
%% to a single short clause each to pay for the larger type; the full enumerations
%% live in the prose of Sections 5 and 6.

\begin{figure}[tbp]
\centering
\ifcbexcalidraw
  \begin{cbwide}\cbdiagram{fig-system}\end{cbwide}
\else
\begin{cbwide}
\begin{tikzpicture}[x=1cm,y=1cm]

  \def\MX{3.30}   % process band left edge
  \def\MW{9.90}   % process band width
  \def\RX{13.66}  % state column left edge
  \def\RW{4.06}   % state column width
  \def\GA{13.36}  % lane A: control plane out to MLflow
  \def\GB{17.70}  % lane B: the production agent's round trip
  \def\CH{-7.72}  % channel below the band for long horizontal edges

  % ======================= clients ==========================================
  \node[anchor=north west, font=\figB\bfseries, text=cbMuted] at (0,0.46) {CLIENTS};
  \node[surface, anchor=north west, text width=2.44cm, minimum height=0.68cm]
    (br) at (0,-0.10) {\textbf{Browser}\\{\figB React SPA session}};
  \node[surface, anchor=north west, text width=2.44cm, minimum height=0.68cm]
    (ap) at (0,-0.98) {\textbf{API client}\\{\figB account-scoped}};
  \node[surface, anchor=north west, text width=2.44cm, minimum height=0.68cm]
    (sv) at (0,-1.86) {\textbf{Service caller}\\{\figB published workflow}};
  \node[extern, anchor=north west, text width=2.44cm, minimum height=0.92cm]
    (pa) at (0,-2.84)
    {\textbf{Production agent}\\{\figB the consumer, outside\\the control plane}};

  % ======================= the ASGI process =================================
  \node[band, anchor=north west, minimum width=\MW cm, minimum height=7.34cm]
    (proc) at (\MX,0.16) {};
  \node[bandlab, anchor=north west, text width=9.6cm, align=left] at (\MX+0.12,0.50)
    {ONE \sysbare{} ASGI APPLICATION \normalfont\itshape (embedded in \mlflow{}
     or standalone --- \figref{fig:topologies})};

  % ---- admission -----------------------------------------------------------
  \node[govern, anchor=north west, text width=\MW cm-0.60cm, minimum height=0.56cm]
    (adm) at (\MX+0.24,-0.10)
    {\textbf{Admission}\ \ {\figB session or token \dt scopes \dt CSRF \dt rate
     limit}};

  % ---- surfaces ------------------------------------------------------------
  \foreach [count=\i from 0] \t in {%
      {\textbf{Route handlers}\\{\figB \statRouteDecls{} declarations}},
      {\textbf{SSE + webhooks}\\{\figB live events}},
      {\textbf{\aria tool loop}\\{\figB risk-tiered}}}
  {
    \pgfmathsetmacro{\w}{(\MW-0.48-2*0.18)/3}
    \node[surface, anchor=north west, text width=\w cm-8pt, minimum height=0.78cm,
          inner xsep=2pt] at ({\MX+0.24+\i*(\w+0.18)}, -0.80) {\t};
  }

  % ---- lifecycle modes -----------------------------------------------------
  \node[control, anchor=north west, text width=\MW cm-0.60cm, minimum height=0.50cm,
        font=\figB] (modes) at (\MX+0.24,-1.72)
    {\mode{Author}\arrowto\mode{Test}\arrowto\mode{Evaluate}\arrowto%
     \mode{Calibrate}\arrowto\mode{Release}\arrowto\mode{Observe}};

  % ---- control-plane services ----------------------------------------------
  \foreach [count=\i from 0] \t in {%
      {\textbf{Orchestrator}\\{\figB triage, diagnose,\\candidate}},
      {\textbf{Evaluation}\\{\figB scorecards,\\LLM judges}},
      {\textbf{Calibration}\\{\figB \statOptimizers{} optimizer\\paths, replay}},
      {\textbf{Apply +}\\\textbf{promote}\\{\figB gate, alias,\\checkpoint}}}
  {
    \pgfmathsetmacro{\w}{(\MW-0.48-3*0.18)/4}
    \node[control, anchor=north west, text width=\w cm-8pt, minimum height=1.30cm,
          inner xsep=2pt] at ({\MX+0.24+\i*(\w+0.18)}, -2.36) {\t};
  }

  % ---- governance substrate ------------------------------------------------
  \node[govern, anchor=north west, text width=\MW cm-0.60cm, minimum height=0.48cm,
        font=\figB] (gov) at (\MX+0.24,-3.82)
    {execution policy \dt gate verdicts \dt release control \dt audit log \dt
     effect ledger};

  % ---- kernel --------------------------------------------------------------
  \node[control, anchor=north west, text width=\MW cm-0.60cm, minimum height=0.48cm,
        font=\figB] (kern) at (\MX+0.24,-4.44)
    {config \& secrets \dt persistence \dt storage \dt tool execution \dt event
     bus \dt provider adapters};

  % ---- durable queues + loops ---------------------------------------------
  \node[db, anchor=north west, text width=2.20cm, minimum height=1.86cm]
    (q) at (\MX+0.24,-5.10)
    {\textbf{Durable}\\\textbf{queues}\\[0.4ex]{\figB status, claim\\and lease
      columns\\on rows}};
  \node[async, anchor=north west, text width=\MW cm-3.60cm, minimum height=1.86cm]
    (loops) at (\MX+0.24+2.86,-5.10)
    {\textbf{Up to \statLoops{} in-process loops}\\{\figB no separate worker tier}\\[0.4ex]
     {\figB Refinement \dt CalibrationDrain \dt WorkflowRun\\
      AriaPlan \dt KnowledgeBase \dt Scheduler\\
      Janitor \dt WebhookDispatcher}};
  \draw[flow, draw=cbAsync] ($(q.east)+(0,0.24)$) -- ($(loops.west)+(0,0.24)$);
  \draw[flow, draw=cbAsync] ($(loops.west)-(0,0.24)$) -- ($(q.east)-(0,0.24)$);

  % ======================= state and external systems =======================
  \node[anchor=north west, font=\figB\bfseries, text=cbMuted] at (\RX,0.46)
    {STATE \& EXTERNAL};
  \node[extern, anchor=north west, text width=\RW cm-8pt, minimum height=0.94cm]
    (mlf) at (\RX,-0.10)
    {\textbf{\mlflow{} 3.14+}\\{\figB prompt versions + aliases\\traces \dt
     assessments}};
  \node[db, anchor=north west, text width=\RW cm-8pt, minimum height=0.94cm]
    (pg) at (\RX,-1.24)
    {\textbf{PostgreSQL 17}\\{\figB \statModels{} tables --- metadata\\\emph{and}
     the file inventory}};
  \node[db, anchor=north west, text width=\RW cm-8pt, minimum height=0.76cm]
    (obj) at (\RX,-2.38)
    {\textbf{Object storage}\\{\figB file bytes only}};
  \node[extern, anchor=north west, text width=\RW cm-8pt, minimum height=0.76cm]
    (llm) at (\RX,-3.34)
    {\textbf{LLM providers}\\{\figB OpenAI \dt Claude \dt Gateway}};
  \node[extern, anchor=north west, text width=\RW cm-8pt, minimum height=0.76cm]
    (mcp) at (\RX,-4.30)
    {\textbf{MCP servers}\\{\figB containment or sidecar}};
  \node[extern, anchor=north west, text width=\RW cm-8pt, minimum height=0.76cm]
    (tx) at (\RX,-5.26)
    {\textbf{Event transport}\\{\figB NATS \dt Redis \dt DB}};

  % ======================= edges ============================================
  % Account clients and service callers all funnel through admission.
  \draw[flow] (br.east) -- ({\MX+0.22},-0.32);
  \draw[flow] (ap.east) -- ({\MX+0.22},-0.38);
  \draw[flow] (sv.east) -- ({\MX+0.22},-0.44);

  % The layer spine: each row hands off to the one below it.
  \foreach \ya/\yb in {-0.66/-0.76, -1.58/-1.68, -2.22/-2.32,
                       -3.66/-3.78, -4.30/-4.40, -4.92/-5.06}
    { \draw[flow, draw=cbMuted] ({\MX+\MW/2},\ya) -- ({\MX+\MW/2},\yb); }

  % Only the kernel touches state; lane A carries the MLflow-owned effects.
  \draw[flow] (kern.east) -- (pg.west);
  \draw[flow] (kern.east) -- (obj.west);
  \draw[flow] (kern.east) -- (llm.west);
  \draw[flow] (kern.east) -- (mcp.west);
  \draw[flow] (loops.east) -- (tx.west);
  \draw[flow] ({\MX+\MW-0.12},-2.90) -- (\GA,-2.90) -- (\GA,-0.44) -- (mlf.west);

  % Lane B: the round trip that closes the loop, kept outside both the process
  % band and the state column so that it crosses nothing.
  \draw[{Stealth[length=4pt,width=3pt]}-{Stealth[length=4pt,width=3pt]},
        draw=cbExtern, line width=0.8pt]
    (pa.south) -- (1.22,\CH) -- (\GB,\CH) -- (\GB,-0.60) -- (mlf.east);
  \node[anchor=south, font=\figB, text=cbInk, align=center, inner sep=1.8pt]
    at (9.6,\CH+0.04)
    {the agent resolves \prodalias{} at call time and emits traces --- when
     \mode{Release} rotates the alias, the next call runs the new version with
     \textbf{no client redeploy}};

\end{tikzpicture}
\end{cbwide}
\fi
\caption{High-level architecture and component interaction. Every account client
and service caller enters through one admission stage; \Lmode{} drives four
control-plane service groups; and the kernel is the only tier that touches state
or external systems. Two properties are load-bearing for the rest of the paper.
First, \sys is a \emph{sibling} surface to \mlflow{}, not a proxy in front of it:
each owns its own store, which is what makes the reconciliation boundary in
\figref{fig:dataflow} necessary. Second, there is no separate worker tier ---
queued work is arbitrated through claim and lease columns on relational rows, so
one process both serves requests and drains queues. The paired arrows between the
queue and the loop group are a consumer's atomic claim and the status transition it
writes back (\algref{alg:claim}). The
outer lane is the only path that leaves the platform entirely, and it is why a
release needs no client deployment.}
\label{fig:system}
\figkey
\end{figure}


\sys is one ASGI application, plus a React single-page application it serves from
the same origin, plus the state it owns. \figref{fig:layers} gives the layered
view and \figref{fig:system} the runtime view; this section introduces the three
concepts a reader needs before either diagram is useful.

\subsection{The governed asset}

%% fig-asset.tex -- Figure (single column): the twelve facets of a governed asset.
%%
%% This is the unit of value. The figure names twelve *possible* facets; the
%% shading distinguishes those that are structurally universal from those a
%% family implements only if its idiom supports them. Table tab:families is the
%% row-by-row truth; this figure is the vocabulary that table uses.
%%
%% Layout note. This is a uniform 3x4 grid, so every cell is the same height and
%% that height is stated once, computed from the worst case: one title line at
%% \figH plus two detail lines at \figB, plus the card's inner padding. Every
%% detail string is therefore held to what fits two lines at the cell's text
%% width -- roughly forty characters. An earlier version let the details run to
%% whatever length read well and pinned the row pitch by eye, which put text
%% through the border of the cell below it. Nothing in LaTeX warns about that: the
%% page is not overfull, only the box is.

\begin{figure}[tbp]
\centering
\begin{cbscaled}[0.94]
\begin{tikzpicture}[x=1cm,y=1cm]

  \def\W{8.10}
  \def\G{0.18}
  \def\CH{1.16}   % cell height: title + two detail lines + padding
  \def\CP{1.32}   % row pitch

  \foreach [count=\i from 0] \s/\n/\d in {%
      asset/{Typed definition}/{schema-validated spec},
      asset/{Version history}/{immutable version\\snapshots},
      muted/{Live target}/{alias selecting what\\runs now},
      muted/{Test surface}/{runs with recorded\\inputs and outputs},
      muted/{Evidence base}/{test sets, traces,\\assessments},
      muted/{Evaluation}/{scorecards, scorers,\\authored judges},
      muted/{Gate semantics}/{advancement, verdict,\\deploy},
      muted/{Calibration}/{proposes a better\\version},
      muted/{Release path}/{audited promote\\and rollback},
      asset/{Authority}/{scopes and risk tiers\\on every mutation},
      asset/{Trace \& audit}/{traces, audit rows,\\incidents},
      muted/{Packaging}/{manifests, exports,\\service exposure}}
  {
    \pgfmathsetmacro{\col}{int(mod(\i,3))}
    \pgfmathsetmacro{\row}{int(floor(\i/3))}
    \pgfmathsetmacro{\w}{(\W-2*\G)/3}
    \node[\s, anchor=north west, text width=\w cm-9pt, minimum height=\CH cm,
          font=\figB, align=left, inner xsep=3pt]
      at ({\col*(\w+\G)}, {-\row*\CP})
      {\textbf{\figH\n}\\[0.35ex]\d};
  }

  % ---- the legend that carries the argument --------------------------------
  % Placed from the grid's own extent rather than at a guessed y.
  \node[anchor=north west, font=\figB, align=left, text width=\W cm, inner sep=0pt]
    at (0, {-3*\CP-\CH-0.16})
    {\keyswatch{asset}{}~\,present for every family\qquad
     \keyswatch{muted}{}~\,family-specific (\tabref{tab:families})};

\end{tikzpicture}
\end{cbscaled}
\caption{Anatomy of a \gasset: the unit of value the whole platform exists to
produce. Twelve facets are \emph{possible}; a family implements and explicitly
wires only those its idiom supports. The four solid facets hold for every family
because they are what makes a record governed at all --- a typed definition,
version history, an authority model, and a trace and audit trail. The eight muted
facets do not transfer: test sets and judges are evidence and scoring assets with
no \livetarget{} and no release path, tools expose family history without a live
alias, and gate semantics differ per family rather than being one cross-artifact
policy. What a \gasset is \emph{not}, stated plainly: not a prompt file in
version control, not a dashboard over traces, and not an agent framework. It is a
typed record carrying enough version, evidence, and release metadata that a change
to it is reviewable after the fact.}
\label{fig:asset}
\end{figure}


The unit of value is the \gasset{}: a typed record that carries enough metadata
that a change to it is reviewable after the fact. \figref{fig:asset} enumerates
twelve possible facets. Four hold for every family, because they are what makes a
record governed at all:

\begin{itemize}
\item a \textbf{typed definition} --- a schema-validated specification that is the
      source of truth for the asset, not a rendering of it;
\item \textbf{version history} --- immutable snapshots, or immutable registry
      versions held externally;
\item an \textbf{authority model} --- which of the \statScopesW{} RBAC scopes may
      read, mutate, and release it;
\item a \textbf{trace and audit trail} --- the durable record of what was done to
      it and by whom.
\end{itemize}

The remaining eight --- live target, test surface, evidence base, evaluation, gate
semantics, calibration, release path, packaging --- are \emph{family-specific}, and
that asymmetry is the subject of \secref{sec:arch}.

It is worth being explicit about what a \gasset is not. It is not a prompt file
under version control: version control gives history but no live target, no bound
evidence, and no authority model. It is not a dashboard over traces: a dashboard
describes behaviour but owns no artifact. And it is not an agent: agents are
composed in a framework of the team's choosing and are \emph{clients} of governed
assets.

\subsection{Six lifecycle modes}

\statModesC{} verbs act on those nouns, and they are deliberately few:
\mAuthor{} (draft, render, compile, validate), \mTest{} (bounded runs against
fixtures), \mEval{} (scorecards, deterministic scorers, authored LLM judges,
review queues), \mCal{} (propose a measurably better version), \mRel{} (gate,
promote, roll back), and \mObs{} (trace, meter, incident, replay).

The modes are reusable concepts, not a state machine every asset traverses. A test
set is authored, tested against, and observed, but is never released --- it is the
evidence others are released against. A judge is authored and evaluated for human
agreement, but has no live target. Enumerating the modes separately from the
families is what makes the coverage matrix in \tabref{tab:families} expressible at
all: without that separation there is only a single undifferentiated notion of
``supported.''

\subsection{The governance chain}

%% fig-chain.tex -- Figure: the governance chain in three registers.
%%
%% Upper: the seven-term abstract chain -- a reference shape, NOT a stage count.
%% Middle: the concrete six-stage prompt refinement path, the fullest instance.
%% Lower: the durable residue each term leaves behind.
%%
%% Register labels are rotated into the left margin so the vertical space between
%% registers stays clear for the correspondence lines, which are deliberately not
%% one-to-one -- that mismatch is the figure's argument.
%%
%% Type scale: \figH headings, \figB detail. Seven columns across 7in leaves each
%% cell about 2.1cm, so detail lines are held to two or three short words.

\begin{figure}[tbp]
\centering
\begin{cbwide}
\begin{tikzpicture}[x=1cm,y=1cm]

  \def\X{1.86}    % content left edge (left of it: rotated register labels)
  \def\W{15.60}   % content width
  \def\G{0.20}

  % ================= upper register: the abstract chain =====================
  \foreach [count=\i from 0] \term/\leaves in {%
      {\mode{Signal}}/{a failure worth\\acting on},
      {\mode{Evidence}}/{traces + test\\sets assembled},
      {\mode{Candidate}}/{a proposed\\better version},
      {\mode{Measure-}\\\mode{ment}}/{scored on that\\evidence},
      {\mode{Decision}}/{an operator\\applies, or not},
      {\mode{Release}}/{the \livetarget{}\\moves},
      {\mode{Trace}}/{the next signal\\arrives}}
  {
    \pgfmathsetmacro{\w}{(\W-6*\G)/7}
    \node[control, anchor=north west, text width=\w cm-8pt, minimum height=1.18cm,
          inner xsep=2pt] (c\i) at ({\X+\i*(\w+\G)}, 0.30)
      {\term\\[0.25ex]{\figB\leaves}};
  }
  \foreach \i [evaluate=\i as \j using int(\i+1)] in {0,...,5}
    { \draw[ctlflow] (c\i.east) -- (c\j.west); }

  % the loop closes, routed above the register so the lane below stays clear
  \draw[feedback] (c6.north) -- ({\X+\W-1.0},0.78) -- (2.96,0.78) -- (c0.north);
  \node[anchor=south, font=\figB, text=cbAsync, inner sep=1.8pt]
    at ({\X+\W/2},0.80)
    {the loop closes --- client code keeps resolving \prodalias{}, so the next
     production trace \emph{is} the next signal};

  % ================= middle register: the concrete path =====================
  \foreach [count=\i from 0] \s/\name/\detail in {%
      muted/{Production trace}/{\mlflow{} trace\\+ assessment},
      human/{1.~Verify}/{\textbf{human} ---\\is it real?},
      control/{2.~Diagnose}/{LLM root cause\\on the evidence},
      control/{3.~Optimize}/{policy-selected\\optimizer},
      control/{4.~Evaluate}/{EvalProvider +\\regression gate},
      human/{5.~Apply}/{\textbf{operator} ---\\diff + scores},
      extern/{6.~Promote}/{audited alias\\rotation}}
  {
    \pgfmathsetmacro{\w}{(\W-6*\G)/7}
    \node[\s, anchor=north west, text width=\w cm-8pt, minimum height=1.18cm,
          inner xsep=2pt] (s\i) at ({\X+\i*(\w+\G)}, -1.96)
      {\textbf{\name}\\[0.25ex]{\figB\detail}};
  }
  \foreach \i [evaluate=\i as \j using int(\i+1)] in {0,...,5}
    { \draw[flow] (s\i.east) -- (s\j.west); }

  % ---- the correspondences, drawn because they are not one-to-one ----------
  \foreach \a/\b in {0/1, 1/2, 2/3, 3/4, 4/5, 5/6, 0/0, 1/1}
    { \draw[weak, densely dotted, draw=cbControl, line width=0.5pt]
        (c\a.south) -- (s\b.north); }

  % ================= lower register: the durable residue ====================
  \foreach [count=\i from 0] \r in {%
      {verification\\item},
      {refinement job\\+ evidence},
      {diagnosis +\\candidate},
      {scores + gate\\decision},
      {apply action +\\provenance anchor},
      {checkpoint +\\audit row},
      {new traces,\\SLO incidents}}
  {
    \pgfmathsetmacro{\w}{(\W-6*\G)/7}
    \node[store, anchor=north west, text width=\w cm-8pt, minimum height=0.86cm,
          font=\figB, inner xsep=2pt] at ({\X+\i*(\w+\G)}, -3.56) {\r};
    \draw[flow, draw=cbStore]
      ({\X+\i*(\w+\G)+\w/2}, -3.18) -- ({\X+\i*(\w+\G)+\w/2}, -3.52);
  }

  % ================= rotated register labels ================================
  % Short labels on purpose: a rotated multi-word label is taller than the register
  % it names, so it runs into the neighbouring one.
  \foreach \y/\lab in {-0.29/{CONCEPTS}, -2.55/{STAGES}, -3.99/{RECORDS}}
  {
    \node[anchor=center, font=\figB\bfseries, text=cbMuted] at (0.74,\y)
      {\rotatebox{90}{\lab}};
  }

\end{tikzpicture}
\end{cbwide}
\caption{The \chainname{} in three registers. The top row is a \emph{reference
shape}, not a worker pipeline: it names seven concepts a refinement-capable
\family may occupy. The middle row is the concrete prompt path --- the fullest
realization of that shape --- with six numbered stages and exactly two human
decisions. The dotted correspondences are deliberately not one-to-one, and the
mismatch is the figure's argument: \cEvidence{} has no stage of its own, because
evidence assembly is folded into the transition from \mode{Verify} toward
\mode{Diagnose}; \cSignal{} spans both an incoming production trace and the human
confirmation that the failure is real; and \cTrace{} is simultaneously the
seventh term and the input to the next iteration. Reading the seven terms as
seven worker stages is the single most common misreading of this architecture.
Other families reuse the shape without inheriting its guarantees: a workflow is
measured by manifest replay, a tool by a revision-fenced deterministic suite, a
knowledge base by retrieval calibration (\tabref{tab:families}). The bottom row
is what separates this from an observability dashboard --- every term deposits
durable state, so a release is reconstructable from records rather than from
recollection.}
\label{fig:chain}
\figkey
\end{figure}


The chain is what makes a change reviewable, and \figref{fig:chain} presents it in
three registers. The upper register names \statChainTermsW{} concepts:
\cSignal{} (a failure worth acting on), \cEvidence{} (traces and examples
assembled into a corpus), \cCandidate{} (a proposed better version), \cMeasure{}
(scored against that corpus), \cDecision{} (an operator applies, or does not),
\cRelease{} (the live target moves), and \cTrace{} (the next signal arrives).

We stress a point the figure is drawn to make and that is routinely misread: these
are \emph{seven concepts, not seven stages}. The concrete prompt path has six
numbered stages, and the correspondence to the seven terms is not bijective.
\cEvidence{} has no stage of its own --- evidence assembly is folded into the
transition from \mode{Verify} toward \mode{Diagnose}. \cSignal{} spans both an
incoming production trace and the human confirmation that the failure is real.
\cTrace{} is simultaneously the seventh term and the input to the next iteration.
A reader who counts stages and finds six where the text says seven has found this
non-bijection, not an inconsistency.

The lower register of \figref{fig:chain} is the part that distinguishes this
architecture from an observability pipeline. Each term deposits durable state: a
verification item, a refinement job with its assembled evidence, a diagnosis and
candidate artifact, scores with an enforced gate decision, an explicit apply action
with a provenance anchor, a rollback checkpoint with an audit row, and new traces.
Six weeks later, the question ``why is the prompt like this, and what did we know
when we changed it'' is answerable from records rather than from recollection.
That property --- not any individual mechanism --- is what \sys exists to provide.

\subsection{Reading the two architectural diagrams}

\figref{fig:layers} reads bottom-up. \Linfra{} is the substrate \sys runs on, all
of it swappable by configuration. \Lkernel{} holds the modular services built at
application startup. \Lgov{} is the substrate of shared primitives that asset
paths wire explicitly. \Lfam{} holds the \statFamiliesW{} governed families --- the
nouns. \Lmode{} holds the verbs. \Lsurf{} exposes the result. The dashed left rail
holds platform services --- the evidence base, evaluation, calibration, the
capability registry, the integration hub, project scoping --- that cut across all
six layers rather than occupying one.

\figref{fig:system} reads left-to-right and shows the runtime instead. Every
account client and service caller passes through one admission stage; lifecycle
modes drive four control-plane service groups; the kernel is the only tier that
touches state or external systems; and queued work is drained by up to
\statLoopsW{} in-process loops arbitrated through claim and lease columns on
relational rows. Two properties are load-bearing for everything that follows.
First, \sys is a \emph{sibling} surface to \mlflow{} rather than a proxy in front
of it --- each owns its own store, which is why the reconciliation boundary
described in \secref{sec:exec} exists and cannot be designed away. Second, there is no
separate worker tier, which is a deliberate trade with consequences we spell out
in \secref{sec:exec}.
